% This file is part of JTS-MM (https://github.com/paulotto/jts-mm).
% SPDX-FileCopyrightText: 2026 Paul-Otto Müller
% SPDX-License-Identifier: CC-BY-4.0

\section{The concepts in detail}\label{sec:concepts}
Each concept is shown on the right side of the arch: an overall view, the first molar from the
buccal and from the lingual side, and a section through the first molar. The section is an
editable SVG in \texttt{figures/svg/}; the geometry is in \texttt{stl/} in the coordinate frame of
\texttt{Teeth\_02}, so it drops onto the teeth when imported into the FreeCAD document. The STLs
are concept geometry on a \SI{0.2}{mm} voxel grid, meant as a starting point for CAD, not as
print files.

\newcommand{\conceptfig}[2]{%
\begin{figure}[H]\centering
\begin{subfigure}[b]{0.54\textwidth}\includegraphics[width=\linewidth]{render_#1_obl}\caption{overview}\end{subfigure}\hfill
\begin{subfigure}[b]{0.44\textwidth}\includegraphics[width=\linewidth]{section_#1}\caption{section through the first molar}\end{subfigure}\\[2mm]
\begin{subfigure}[b]{0.49\textwidth}\includegraphics[width=\linewidth,height=45mm,keepaspectratio]{render_#1_molB}\caption{first molar, buccal side}\end{subfigure}\hfill
\begin{subfigure}[b]{0.49\textwidth}\includegraphics[width=\linewidth,height=45mm,keepaspectratio]{render_#1_molL}\caption{first molar, lingual side}\end{subfigure}
\caption{#2}\label{fig:#1}\end{figure}}

% ---------------------------------------------------------------------------
\subsection{Concept 1: snap shell}
\conceptfig{c1}{Concept 1. One conforming shell: lingual flange on all teeth, buccal reciprocal arms on both molars, bridges over the embrasures, labial bar with the marker stem.}
\textbf{Principle.} Like a partial clear aligner: the lingual flange snaps under the lingual height
of contour of the molars; the wall itself is the spring.

\textbf{Parameters.} Wall \SI{1.5}{mm} (PETG/PA12) or \SIrange{1.0}{1.2}{mm} (splint resin);
flange from \SI{1.4}{mm} below the lingual cusp tips to \SIrange{1}{1.5}{mm} onto the lingual
gingiva; buccal arms \SI{1.3}{mm}, above the buccal height of contour; embrasure bridges
\SI{\ge1.5}{mm} wide, \SI{1.0}{mm} thick; labial bar \SI{1.6}{mm} on the cervical half of the
incisors and canines.

\textbf{Build.}
\begin{enumerate}
\item Survey and block out the molar lingual undercut to \SI{0.3}{mm} (PETG) or
\SI{0.5}{mm} (flexible splint resin); relieve all other surfaces by \SI{0.1}{mm}.
\item Offset the blocked-out model by the wall thickness (e.g.\ Meshmixer ``offset'', Blender
``solidify'', or start from \texttt{stl/c1\_printed.stl}) and trim to the outline in \cref{fig:c1}.
\item Print in splint resin (SLA/DLP). In FDM, orient the part so that the flange bends
\emph{within} the layer plane, not across layers.
\item Fit on the printed model, then reline the lingual flange with resin through retention holes.
\end{enumerate}
\textbf{Watch out.} Spring and reference are the same part: whenever the flange flexes, the marker
moves. Preload depends directly on print accuracy.

% ---------------------------------------------------------------------------
\subsection{Concept 2: clothespin jaw}
\conceptfig{c2}{Concept 2. Rigid frame (lingual bar, buccal jaw, bridge over the PM2/M1 embrasure, labial bar) and a compliant lingual jaw hung on a C-loop flexure, an intraoral version of a printed spring clamp.}
\textbf{Principle.} The lingual jaw is carried by a C-shaped loop spring that sits next to the
tongue, below the occlusal plane. It is printed over-closed; inserting the clamp opens it over the
height of contour, and it springs back into the undercut.

\textbf{Parameters.} Loop radius $R=\SI{3.0}{mm}$ to the wall centre, wall $t=\SI{0.8}{mm}$ (two
\SI{0.4}{mm} extrusion lines), width $w=\SI{5.5}{mm}$ along the arch, \SI{270}{\degree} arc,
loop centre at the contact height ($e=0$), over-closure \SI{0.2}{mm}; lingual jaw \SI{1.2}{mm} from
$h=6.3$ to $9.3$; rigid buccal jaw \SI{1.4}{mm}. This gives \SI{\approx1.5}{N} preload and
\SI{\approx3.7}{N} while passing a \SI{0.3}{mm} undercut, at \SI{\approx1.0}{\percent} peak strain
(\cref{sec:loop}). Dimensioned drawings: \cref{fig:drw_c2}.

\textbf{Build.}
\begin{enumerate}
\item Design the loop module (jaw + loop + foot) as a separate clip-in part with a dovetail on the
lingual bar. It can then be printed flat with the loop in the layer plane, and swapped to change
the force.
\item Print the frame rigid (PETG, \SI{100}{\percent} infill, or splint resin).
\item Measure the jaw force with a spring scale at \SI{0.2}{mm} opening and at \SI{0.5}{mm}. Adjust
the width $w$ (force $\propto w$), not the wall thickness: $t$ sets the strain, $w$ sets the force.
\item Add a small tab on the lingual jaw so it can be pulled open with a probe for removal.
\end{enumerate}
\textbf{Watch out.} Printed polymer flexures relax at \SI{37}{\celsius} under constant strain:
measure the force again after \SI{30}{min} in warm water. A loop of $R=\SI{3}{mm}$ reaches about
\SI{7}{mm} lingual of the tooth, into the tongue space; the two-loop variant in \cref{tab:loop}
is more compact.

% ---------------------------------------------------------------------------
\subsection{Concept 3: wire clasps}
\conceptfig{c3}{Concept 3. Rigid printed lingual plate with occlusal rests and an Adams clasp of \SI{0.7}{mm} spring-hard stainless steel around each first molar.}
\textbf{Principle.} The classic removable-retainer solution. The two arrowheads of each Adams
clasp engage the mesio- and disto-buccal undercuts below the contact points, which exist even where
the buccal face itself has only \SI{0.1}{mm}. The spring is metal, so the preload does not creep.

\textbf{Parameters.} \SI{0.7}{mm} hard stainless on molars (\SI{0.6}{mm} on premolars); arrowheads
\SIrange{0.5}{1}{mm} into the undercut; bridge about \SI{1}{mm} off the buccal surface at mid-crown
height; tags cross the embrasures resting on the marginal ridges and end in the plate with
retention bends.

\textbf{Build.}
\begin{enumerate}
\item On a printed or stone model, trim a small notch (about \SI{0.5}{mm}) at the two arrowhead
sites just below the contact points; this is where the arrowheads will sit.
\item Bend the clasp (a dental technician does this in minutes; any orthodontic lab textbook
shows the sequence). Alternative without bending: prefabricated ball clasps in the buccal
embrasures.
\item Print the lingual plate with channels for the tags; fix the tags with cold-cure resin.
\item Adjust retention with pliers at the arrowheads, never by heating the wire.
\end{enumerate}
\textbf{Watch out.} The wire crosses the occlusal embrasures in four places: check the bite.

% ---------------------------------------------------------------------------
\subsection{Concept 4: spring arch (inside push)}
\conceptfig{c4}{Concept 4. A \SI{0.9}{mm} stainless arch behind the teeth (grey: mounted, orange: free shape, pre-spread \SI{0.7}{mm} per side) pushes resin-relined pads outward under the lingual heights of contour of both molars. Rigid anterior yoke and labial bar carry the marker.}
\textbf{Principle.} The expander principle, turned into a clamp: one U-shaped arch, bent slightly
too wide, pushes both molar pads outward into the deepest undercuts in the mouth. The pads seat
themselves (\cref{fig:mech}b). The front of the arch is bridged by the rigid yoke; only the two
legs (\SI{\approx18}{mm} from canine to first molar) act as springs.

\textbf{Parameters.} \SI{0.9}{mm} spring-hard stainless; pre-spread \SI{0.7}{mm} per side; lingual
undercut blocked out to \SI{0.4}{mm}; result \SI{\approx2.3}{N} preload per side, peak stress
\SI{\approx0.9}{GPa} (\cref{sec:springs}). Optional omega loop in each leg for more margin and
adjustability.

\textbf{Build.}
\begin{enumerate}
\item Block out the molar lingual undercut to \SI{0.4}{mm} on the model.
\item Bend the arch on the model, \SI{1}{mm} off the lingual gingiva, from the first-molar pad on
one side to the other; finish the leg ends as small retention loops.
\item Print the pads (splint resin or PETG) with a channel for the loop; fix with cold-cure resin;
reline the tooth side.
\item Activate: open each leg so the distance between the pad contact points is
\SI{1.4}{mm} larger than on the model. Check with a spring scale: \SI{\approx2}{N} per side at
\SI{0.7}{mm}.
\item Remove by pinching the two legs toward each other.
\end{enumerate}
\textbf{Watch out.} Do not pre-spread more than about \SI{1}{mm} or leave deeper undercuts: the
wire then yields (\cref{fig:springs}). Pull-off force is only \SI{\approx1}{N} per side; if that is
not enough, combine with concept~3.

% ---------------------------------------------------------------------------
\subsection{Concept 5: rigid clamp with TPU liner}
\conceptfig{c5}{Concept 5. Hard buccal jaw and hard rest over the PM2/M1 embrasure define the position; a \SIrange{1.5}{2}{mm} TPU liner on the rigid lingual jaw provides the preload.}
\textbf{Principle.} Nothing rigid flexes. The TPU liner is compressed while the clamp passes the
height of contour and expands back into the undercut. TPU is roughly 30--100$\times$ softer than
PETG, so it must be loaded in compression, not used as a clasp arm.

\textbf{Parameters.} Liner \SIrange{1.5}{2}{mm}, TPU 85A--95A, \textbf{ribbed} (e.g.\ four vertical
ribs, about \SI{35}{\percent} solid); interference \SIrange{0.1}{0.2}{mm}; lingual undercut
blocked out to \SI{0.3}{mm}. \Cref{tab:tpu} shows why solid liners are a bad idea: a solid pad
gives tens of newtons for a tenth of a millimetre.

\textbf{Build.}
\begin{enumerate}
\item Print the rigid frame (PETG or splint resin) with a dovetail slot in the lingual jaw.
\item Print the TPU liner as a clip-in insert with the ribs running occluso-gingivally, so the
tooth slides along them on insertion. Dual-material printing with an interlocking pattern also
works.
\item Reline the hard buccal jaw and the rest with resin; do not reline the TPU.
\item Set the interference by the insert thickness; measure pull-off with a spring scale.
\end{enumerate}
\textbf{Watch out.} TPU relaxes under constant compression; re-check the force after a typical
session length. A flexible splint resin is the certified alternative for participants.
